\section{一些基本性质}\label{sec:10}
\subsection{良定义性}\label{sec:10a}
抽象地说，需要考虑以某个集合 $I$ 为指标、彼此等价的范畴所构成的系统。这个概念有以下三个版本，只有中间一个不是显然的。\znote{133}{三个版本的区别在于“复合的相容性”被要求到什么程度：严格系统直接令 $F^{i_2,i_0}$ 等于复合；弱系统只要求任意两个范畴等价，不指定等价的选择，因而没有任何相容性可言；协调系统则给定同构 $T$ 并要求\eref{10.1}交换。后者不显然，是因为两个复合同构之差一般给出一个障碍类（在群作用的情形即 §(10c) 讨论的障碍）；只有当 $T$ 由某个严格自然变换 $S$ 构造出来时（如本段的 $T=L_{K^{i_2}}(R_{H^{i_0}}(S^{i_1}))$），障碍才自动消失。[来源：原书 §(10a)，并见 §(10c) 的\stref{笔记}{10.4}。]}
\begin{definition*}
\begin{itemize}
\item \emph{严格系统}是一族范畴 $A^i$（$i\in I$），连同函子 $F^{i_1,i_0}:A^{i_0}\to A^{i_1}$，满足 $F^{ii}=\Id_{A^i}$ 及 $F^{i_2,i_1}\circ F^{i_1,i_0}=F^{i_2,i_0}$。
\item \emph{协调系统}由范畴 $A^i$、函子 $F^{i_1,i_0}$，以及函子同构 $T^{i_2,i_1,i_0}:F^{i_2,i_1}\circ F^{i_1,i_0}\to F^{i_2,i_0}$ 组成，满足两个条件：首先，$F^{ii}=\Id_{A^i}$；类似地，只要 $i_2=i_1$ 或 $i_1=i_0$，$T^{i_2,i_1,i_0}$ 就是恒等变换。其次，对所有 $(i_3,i_2,i_1,i_0)$，以下图交换：
\SourceEquationDiagram{diagram-10-1}{0.57115323}{10.1}
这里 $L,R$ 分别是作用于自然变换的左复合和右复合；它们是\secref{1e}中讨论的 $\Ai$ 构造的经典对应物。
\item \emph{弱系统}是一族范畴 $A^i$，任意两个都等价，但不指定等价的选择。
\end{itemize}

\end{definition*}
假设有范畴 $A^{\mathrm{tot}}$ 及一族全子范畴 $A^i$，使每个嵌入 $H^i:A^i\to A^{\mathrm{tot}}$ 都是等价。暂时固定某个 $i$，对每个
\SourcePage{142}
$X\in\Ob A^{\mathrm{tot}}$，选择一个与之同构的对象 $K^i(X)\in\Ob A^i$，以及指定的同构 $S^i\in\Hom_{A^{\mathrm{tot}}}(K^i(X),X)$；若 $X$ 已属于 $A^i$，就作显然的选择。很容易将 $K^i$ 赋予函子 $A^{\mathrm{tot}}\to A^i$ 的结构，使 $K^i\circ H^i=\Id_{A^i}$。于是 $S^i$ 成为自然变换 $H^i\circ K^i\to\Id_{A^{\mathrm{tot}}}$，并且 $R_{H^i}(S^i)$、$L_{K^i}(S^i)$ 都是平凡的恒等自然变换。对每个 $i$ 这样做，再令
\[
F^{i_1,i_0}=K^{i_1}\circ H^{i_0},\qquad
T^{i_2,i_1,i_0}=L_{K^{i_2}}\bigl(R_{H^{i_0}}(S^{i_1})\bigr),
\]
就能把各 $A^i$ 组成一个协调范畴系统。

严格系统的概念不作改变就适用于 $\Ai$ 范畴；弱系统也一样，只需将等价换为拟等价。对于协调系统，要求有 $\Ai$ 函子 $\F^{i_1,i_0}:\A^{i_0}\to\A^{i_1}$，满足 $\F^{ii}=\Id_{\A^i}$，并有自然变换的上同调类
\[
[T^{i_2,i_1,i_0}]\in\Hom_{H^0(\operatorname{fun}(\A^{i_0},\A^{i_2}))}
(\F^{i_2,i_1}\circ\F^{i_1,i_0},\F^{i_2,i_0}).
\]
同样，当 $i_2=i_1$ 或 $i_1=i_0$ 时，这个类应为恒等元 $[E_{\F^{i_2,i_0}}]$；而且\eref{10.1}的对应图，采用\secref{1e}所定义的左、右复合函子后，在 $H^0(\operatorname{fun}(\A^{i_0},\A^{i_3}))$ 中交换。

与前面一样，固定一个 $\Ai$ 范畴 $\A^{\mathrm{tot}}$ 及一族与它拟等价的全 $\Ai$ 子范畴 $\A^i$，即可构造协调系统。具体说，以 $\mathcal H^i$ 表示嵌入，先像\stref{定理}{2.9}之前的讨论那样，构造拟等价 $\mathcal K^i:\A^{\mathrm{tot}}\to\A^i$，使 $\mathcal K^i\circ\mathcal H^i=\Id_{\A^i}$。由\stref{引理}{1.7}可知，映射
\begin{equation}\label{eq:10.2}
\begin{aligned}
H(\LG_{\mathcal K^i}):&\ \Hom_{H(\operatorname{fun}(\A^{\mathrm{tot}},\A^{\mathrm{tot}}))}
(\mathcal H^i\circ\mathcal K^i,\Id_{\A^{\mathrm{tot}}})\\
&\longrightarrow\Hom_{H(\operatorname{fun}(\A^{\mathrm{tot}},\A^i))}(\mathcal K^i,\mathcal K^i)
\end{aligned}
\end{equation}
是同构。选择在\eref{10.2}下映到恒等自然变换 $[E_{\mathcal K^i}]$ 的类 $[S^i]$。先考察 $\LG_{\mathcal K^i}\circ\RG_{\mathcal H^i}=\RG_{\mathcal H^i}\circ\LG_{\mathcal K^i}$，再应用一次\stref{引理}{1.7}，可知 $H(\RG_{\mathcal H^i})([S^i])=[E_{\mathcal H^i}]$。其余构造与普通范畴的情形相同。

转向 Floer 理论。回顾，给定带角正合辛流形 $M$，$\F(M)^{\mathrm{pr}}$ 的定义涉及额外选择，见\secref{9j}的汇总。设 $I$ 是所有可能选择的集合，以 $\{\F(M)^{\mathrm{pr},i}\}$ 表示所得 $\Ai$ 范畴族。若转到上同调层次，则 $H\F(M)^{\mathrm{pr},i}$ 构成上述意义下的严格系统。这是 Floer 上同调及其 TQFT 结构的良定义性的结果，\secref{8j}已简要讨论过。

为了处理 $\F(M)^{\mathrm{pr}}$ 本身，方便采用上面描述的形式技巧。定义一个总 Fukaya 范畴 $\F(M)^{\mathrm{pr},\mathrm{tot}}$，其对象是由正合 Lagrangian 子流形 $L$ 与指标 $i\in I$ 组成的二元组 $(L,i)$。态射空间照常取 Floer 复形，对每一对 $((L_0,i_0),(L_1,i_1))$ 分别选择正则 Floer 数据，要求当 $i_0=i_1=i$ 时，数据与将 $(L_0,L_1)$
\SourcePage{143}
视为 $\F(M)^{\mathrm{pr},i}$ 中对象时相同。类似地，通过适当选择条带状端与扰动数据来定义 $d\geq2$ 阶复合，要求当涉及的所有对象都有同一个指标 $i$ 时，选择与 $\F(M)^{\mathrm{pr},i}$ 相同。这稍微修改了标准设置，因为允许 $(d+1)$ 穿孔圆盘上的条带状端依赖于标记 $(L_k,i_k)$。然而，利用\secref{9i}中相容标记族的粘合，很容易适当调整\secref{9g}的一致性条件。结果是，总范畴带有典范的全忠实嵌入
\begin{equation}\label{eq:10.3}
\F(M)^{\mathrm{pr},i}\longrightarrow\F(M)^{\mathrm{pr},\mathrm{tot}}.
\end{equation}
对于任意固定的正合 Lagrangian 子流形 $L$ 及指标 $(i_0,i_1)$，对象 $(L,i_0)$、$(L,i_1)$ 在 $H\F(M)^{\mathrm{pr},\mathrm{tot}}$ 中同构。这是上同调层次的陈述，其证明归结为标准的延拓映射论证。因此，嵌入\eref{10.3}都是拟等价。于是可以利用这一点，把 $\{\F(M)^{\mathrm{pr},i}\}$ 组成 $\Ai$ 范畴的协调系统。严格说，它们是不分次的 $\Ai$ 范畴；不过\stref{笔记}{9.7}适用于这里，也适用于本节后面类似的情形。

此外，只要某个关于 Fukaya 范畴的代数构造具有函子性并保持拟等价，这个结论就同样适用于该构造，例如取导出范畴。\znote{135}{对 $\Ai$ 范畴“取导出范畴”，通常指向 $D(\mathcal A)=H^0(\operatorname{mod}\mathcal A)$，其中 $\operatorname{mod}\mathcal A$ 是 $\mathcal A$ 上右 $\Ai$ 模所成的 $\Ai$ 范畴，$H^0$ 只取零次上同调。欲使它成为三角范畴并使拟等价成为等价，还需再取分裂闭包（Karoubi 包络），即把幂等元的像形式地添进去。Yoneda 嵌入 $X\mapsto\hom_{\mathcal A}(-,X)$ 把 $\mathcal A$ 全忠实地放入 $\operatorname{mod}\mathcal A$，其像在分裂闭包下生成整个导出范畴；分裂闭包的存在性正是“幂等元必分裂”这条万有性质的体现。[来源：nLab "Karoubi envelope"，https://\allowbreak ncatlab.org/\allowbreak nlab/\allowbreak show/\allowbreak Karoubi+\allowbreak envelope 与 nLab "Yoneda embedding"，https://\allowbreak ncatlab.org/\allowbreak nlab/\allowbreak show/\allowbreak Yoneda+\allowbreak embedding（2026-11-15 访问）。]}剩下的问题是：构成协调系统的函子与自然变换本身，在多大程度上是典范的？这留给有兴趣的读者。
\begin{note}\label{st:10.1}
在普通范畴的框架中，有一种反向构造：从范畴 $A^i$ 的协调系统出发，构造一个共同范畴 $A^{\mathrm{tot}}$，以及等价 $A^i\to A^{\mathrm{tot}}$。在 $\Ai$ 背景下，这看来不能成立，原因在于我们的协调系统定义本质上太弱，只停留在上同调函子范畴 $H(\operatorname{fun})$ 的层次。
\end{note}

\subsection{自等价}\label{sec:10b}
这个主题与前一个紧密相关；事实上，可以用群胚统一处理。需要讨论群 $G$ 作用于范畴 $A$ 的含义。
\begin{definition*}
\begin{itemize}
\item \emph{严格 $G$ 作用}是一族函子 $F^g:A\to A$（$g\in G$），满足 $F^e=\Id_A$、$F^{g_2g_1}=F^{g_2}\circ F^{g_1}$。
\item \emph{协调 $G$ 作用}由函子 $F^g$ 及同构 $T^{g_2,g_1}:F^{g_2}\circ F^{g_1}\to F^{g_2g_1}$ 组成，满足 $F^e=\Id_A$；当 $g_2=e$ 或 $g_1=e$ 时，$T^{g_2,g_1}$ 为恒等变换；并且以下图交换：
\SourceEquationDiagram{diagram-10-4}{0.46730719}{10.4}
\SourcePage{144}
\item \emph{弱 $G$ 作用}是一族函子 $F^g$，满足 $F^e$ 同构于 $\Id_A$、$F^{g_2g_1}$ 同构于 $F^{g_2}\circ F^{g_1}$，但不指定这些同构的选择。
\end{itemize}

\end{definition*}
假设范畴 $A^{\mathrm{free}}$ 带有严格 $G$ 作用 $\{\Phi^g\}$，$A$ 是其全子范畴，且嵌入 $H:A\to A^{\mathrm{free}}$ 为等价。像\secref{10a}那样，选择 $K:A^{\mathrm{free}}\to A$ 及函子同构 $S:H\circ K\to\Id_{A^{\mathrm{free}}}$。于是，在 $A$ 上得到诱导的协调 $G$ 作用：
\begin{equation}\label{eq:10.5}
F^g=K\circ\Phi^g\circ H,\qquad
T^{g_2,g_1}=L_{K\circ\Phi^{g_2}}\bigl(R_{\Phi^{g_1}\circ H}(S)\bigr).
\end{equation}
现在应当很清楚，如何将这些定义推广到 $\Ai$ 框架。注意，在此还有比严格作用更基本的另一种可能，即第\sr{sec:6}{6 节}中的朴素群作用；不过它们在当前讨论中不起作用。为说明\eref{10.5}的 $\Ai$ 对应物满足协调关系，论证如下。对应于\eref{10.4}的图的两条复合路径，都具有形式
\[
\LG^1_{\mathcal K\circ\Phi^{g_3}}\bigl(\RG^1_{\Phi^{g_2}\circ\Phi^{g_1}\circ\mathcal H}(\text{某个表达式})\bigr),
\]
括号内的表达式分别为
\begin{equation}\label{eq:10.6}
\begin{aligned}
&\mu^2_{\Q}\bigl(S,\LG^1_{\mathcal H\circ\mathcal K}(\LG^1_{\Phi^{g_2}}(\RG^1_{\Phi^{g_2^{-1}}}(S)))\bigr),\\
&\mu^2_{\Q}\bigl(\LG^1_{\Phi^{g_2}}(\RG^1_{\Phi^{g_2^{-1}}}(S)),
\RG^1_{\Phi^{g_2}\circ\mathcal H\circ\mathcal K\circ\Phi^{g_2^{-1}}}(S)\bigr),
\end{aligned}
\end{equation}
其中 $\Q=\operatorname{fun}(\A^{\mathrm{free}},\A^{\mathrm{free}})$。应用\stref{引理}{1.4}，取
\[
\begin{gathered}
\G_1=\mathcal H\circ\mathcal K,\qquad
\G_2=\Phi^{g_2}\circ\mathcal H\circ\mathcal K\circ\Phi^{g_2^{-1}},\\
T_1=S,\qquad T_2=\LG^1_{\Phi^{g_2}}(\RG^1_{\Phi^{g_2^{-1}}}(S)),
\end{gathered}
\]
可知在 $H(\Q)$ 层次，\eref{10.6}中的两项相同。
\begin{note}\label{st:10.2}
与\stref{笔记}{10.1}所述相应，在普通范畴的框架中存在反向构造，将 $A$ 上的协调群作用变为等价范畴 $A^{\mathrm{strict}}$ 上的严格作用；特殊情形 $G=\Z/2$ 见\secref{14b}。同样，在 $\Ai$ 背景下，这看来不能成立，至少对上面给出的协调性定义如此。
\end{note}
固定一个带角正合辛流形 $M$，令 $G=\operatorname{Aut}(M,\partial M)$ 为其相对于边界的自同构群。每个 $\phi\in G$ 都在整个 Floer TQFT 上诱导典范自同构，特别是在 $H\F(M)^{\mathrm{pr}}$ 上诱导自同构；它们构成严格 $G$ 作用。在上链层次，按如下方式进行。定义一个 Fukaya 型的不分次 $\Ai$ 范畴 $\F(M)^{\mathrm{pr},\mathrm{free}}$，其对象是由正合 Lagrangian 子流形 $L\subset M$ 与 $\gamma\in G$ 组成的二元组 $(L,\gamma)$。为每组有限个对象选择 Floer 数据与扰动数据，借此定义态射空间 $CF^{\mathrm{pr}}((L_0,\gamma_0),(L_1,\gamma_1))$ 及它们之间的复合映射，并要求
\SourcePage{145}
这些选择关于 $G$ 作用 $\phi(L,\gamma)=(\phi(L),\phi\circ\gamma)$ 等变。这并无困难，因为该作用在对象上自由。结果是，$\F(M)^{\mathrm{pr},\mathrm{free}}$ 带有严格 $G$ 作用，由高阶项 $d>1$ 全为零的 $\Ai$ 函子 $\phi_*$ 组成。显然，可将 $\F(M)^{\mathrm{pr}}$ 作为全子范畴嵌入 $\F(M)^{\mathrm{pr},\mathrm{free}}$；因为任意两个 $(L,\gamma_0)$、$(L,\gamma_1)$ 在上同调层次同构，这个嵌入是拟等价。因此，得到 $\F(M)^{\mathrm{pr}}$ 上诱导的协调 $G$ 作用。\znote{137}{“该作用在对象上自由”是这里能选出等变数据的关键：若某个 $\phi\neq\id$ 固定了某个 $L$，则同一组 Floer 数据必须同时满足 $L$ 与 $\phi(L)=L$ 两套边界条件，横截性会因此失效，这正是一般等变横截性问题的困难所在。改用对象 $(L,\gamma)$ 后，$\phi$ 把 $(L,\gamma)$ 送到 $(\phi L,\phi\gamma)$，只要 $\phi\neq\id$ 两者就不同，于是可以在每个轨道上任选数据再用 $G$ 推移。代价是对象集变大：$\F(M)^{\mathrm{pr},\mathrm{free}}$ 的对象数约为 $\F(M)^{\mathrm{pr}}$ 的 $|G|$ 倍。[来源：原书 §(10b) 本段。]}
